\chapter{Algebra \& Number Theory}

\section{Group Theory and Representations}
\textbf{Prerequisite Knowledge}: Sets, mappings, equivalence relations, basic linear algebra (matrices, vector spaces).

\begin{definition}[Group Action]
A group action of a group $G$ on a set $X$ is a function $G \times X \to X$, denoted $(g, x) \mapsto g \cdot x$, such that $e \cdot x = x$ for all $x \in X$ (where $e$ is the identity of $G$), and $g \cdot (h \cdot x) = (gh) \cdot x$ for all $g, h \in G$ and $x \in X$. The orbit of $x \in X$ is $O_x = \{g \cdot x \mid g \in G\}$. The stabilizer of $x \in X$ is $G_x = \{g \in G \mid g \cdot x = x\}$.
\end{definition}

\begin{definition}[Representation]
A linear representation of a group $G$ on a vector space $V$ over a field $F$ is a group homomorphism $\rho \colon G \to GL(V)$, where $GL(V)$ is the general linear group of $V$.
\end{definition}

\begin{theorem}[Sylow Theorems]
Let $G$ be a finite group of order $p^n m$, where $p$ is a prime and $p$ does not divide $m$.
1. $G$ has a Sylow $p$-subgroup (a subgroup of order $p^n$).
2. Any two Sylow $p$-subgroups are conjugate in $G$.
3. The number of Sylow $p$-subgroups, $n_p$, divides $m$ and satisfies $n_p \equiv 1 \pmod p$.
\end{theorem}

\begin{example}
Let $G = S_3$, the symmetric group on 3 elements. The order of $S_3$ is 6. A Sylow 2-subgroup has order 2, and there are 3 such subgroups. A Sylow 3-subgroup has order 3, and there is exactly 1 such subgroup, which must be normal.
\end{example}

\section{Galois Theory and Field Extensions}
\textbf{Prerequisite Knowledge}: Polynomial rings, irreducibility, vector spaces.

\begin{definition}[Splitting Field]
A splitting field of a polynomial $f(x) \in F[x]$ is a field extension $E/F$ such that $f(x)$ factors into linear factors in $E[x]$ and $E$ is generated over $F$ by the roots of $f(x)$.
\end{definition}

\begin{definition}[Galois Group]
For a field extension $E/F$, the Galois group $\mathrm{Gal}(E/F)$ is the group of all field automorphisms of $E$ that fix every element of $F$.
\end{definition}

\begin{theorem}[Fundamental Theorem of Galois Theory]
Let $E/F$ be a finite Galois extension (normal and separable) with Galois group $G = \mathrm{Gal}(E/F)$. There is a bijection between the set of subfields $K$ of $E$ containing $F$ and the set of subgroups $H$ of $G$, given by $K \mapsto \mathrm{Gal}(E/K)$ and $H \mapsto E^H$ (the fixed field of $H$).
\end{theorem}

\begin{example}
Consider $f(x) = x^2 - 2$ over $\mathbb{Q}$. Its roots are $\pm \sqrt{2}$. The splitting field is $\mathbb{Q}(\sqrt{2})$. The Galois group $\mathrm{Gal}(\mathbb{Q}(\sqrt{2})/\mathbb{Q})$ has order 2, consisting of the identity and the automorphism $\sqrt{2} \mapsto -\sqrt{2}$.
\end{example}

\section{Ring and Module Theory}
\textbf{Prerequisite Knowledge}: Basic definition of rings, ideals, homomorphisms.

\begin{definition}[PID and UFD]
A Principal Ideal Domain (PID) is an integral domain where every ideal is generated by a single element. A Unique Factorization Domain (UFD) is an integral domain where every non-zero non-unit can be factored uniquely (up to associates and ordering) into irreducible elements. Every PID is a UFD.
\end{definition}

\begin{definition}[Module]
A module over a ring $R$ (or $R$-module) is an abelian group $M$ equipped with an operation $R \times M \to M$ satisfying axioms analogous to those of a vector space over a field.
\end{definition}

\begin{theorem}[Structure Theorem for Finitely Generated Modules over a PID]
Let $R$ be a PID and $M$ a finitely generated $R$-module. Then $M \cong R^r \oplus R/(a_1) \oplus \dots \oplus R/(a_m)$, where $r \geq 0$ is the rank of $M$, and $a_1, \dots, a_m$ are non-zero non-units in $R$ such that $a_1 \mid a_2 \mid \dots \mid a_m$.
\end{theorem}

\begin{example}
Any finitely generated abelian group is a finitely generated $\mathbb{Z}$-module. Since $\mathbb{Z}$ is a PID, any such group is isomorphic to $\mathbb{Z}^r \oplus \mathbb{Z}/(n_1) \oplus \dots \oplus \mathbb{Z}/(n_k)$.
\end{example}

\section{Commutative Algebra and Algebraic Geometry}
\textbf{Prerequisite Knowledge}: Ideals (prime, maximal), polynomial rings, topological spaces.

\begin{definition}[Coordinate Ring]
For an affine algebraic variety $V \subseteq \mathbb{A}^n(k)$ defined by an ideal $I \subseteq k[x_1, \dots, x_n]$, its coordinate ring is $k[V] = k[x_1, \dots, x_n]/I$.
\end{definition}

\begin{definition}[Integral Closure]
Let $R \subseteq S$ be rings. An element $s \in S$ is integral over $R$ if it is a root of a monic polynomial with coefficients in $R$. The set of all elements in $S$ integral over $R$ is the integral closure of $R$ in $S$.
\end{definition}

\begin{theorem}[Hilbert's Nullstellensatz]
Let $k$ be an algebraically closed field. If $I$ is an ideal in $k[x_1, \dots, x_n]$, then the ideal of the affine variety defined by $I$ is the radical of $I$, i.e., $I(V(I)) = \sqrt{I}$.
\end{theorem}

\begin{example}
Let $I = (x^2 - y^3) \subseteq \mathbb{C}[x,y]$. The coordinate ring $\mathbb{C}[x,y]/(x^2-y^3)$ represents the plane curve with a cusp at the origin. It is not integrally closed.
\end{example}

\section{$p$-adic Analysis and Algebraic Number Theory}
\textbf{Prerequisite Knowledge}: Metric spaces, congruences, absolute values.

\begin{definition}[$p$-adic Valuation]
For a prime $p$, the $p$-adic valuation $v_p(n)$ of a non-zero integer $n$ is the exponent of the highest power of $p$ dividing $n$. The $p$-adic absolute value is $|x|_p = p^{-v_p(x)}$. The field of $p$-adic numbers $\mathbb{Q}_p$ is the completion of $\mathbb{Q}$ with respect to this absolute value.
\end{definition}

\begin{definition}[Ring of Integers]
The ring of integers $\mathcal{O}_K$ of an algebraic number field $K$ (a finite extension of $\mathbb{Q}$) is the integral closure of $\mathbb{Z}$ in $K$.
\end{definition}

\begin{theorem}[Hensel's Lemma]
Let $f(x) \in \mathbb{Z}_p[x]$. If there exists an $a \in \mathbb{Z}_p$ such that $f(a) \equiv 0 \pmod p$ and $f'(a) \not\equiv 0 \pmod p$, then there exists a unique $\alpha \in \mathbb{Z}_p$ such that $f(\alpha) = 0$ and $\alpha \equiv a \pmod p$.
\end{theorem}

\begin{example}
Consider $x^2 - 7 = 0$ in $\mathbb{Z}_3$. Since $x^2 \equiv 1 \pmod 3$ for $x=1$ and $2(1) \not\equiv 0 \pmod 3$, Hensel's lemma ensures that 7 has a square root in $\mathbb{Z}_3$.
\end{example}

\section{Linear and Multilinear Algebra}
\textbf{Prerequisite Knowledge}: Vector spaces, linear transformations, matrices, determinants.

\begin{definition}[Tensor Product]
Let $V$ and $W$ be vector spaces over a field $F$. Their tensor product $V \otimes W$ is a vector space together with a bilinear map $V \times W \to V \otimes W$ that satisfies a universal property for all bilinear maps from $V \times W$.
\end{definition}

\begin{definition}[Nilpotent Operator]
A linear operator $T \colon V \to V$ is nilpotent if $T^k = 0$ for some positive integer $k$.
\end{definition}

\begin{theorem}[Cayley--Hamilton Theorem]
Every square matrix over a commutative ring satisfies its own characteristic equation. If $p(\lambda) = \det(\lambda I - A)$ is the characteristic polynomial of an $n \times n$ matrix $A$, then $p(A) = 0$.
\end{theorem}

\begin{example}
Let $A = \begin{pmatrix} 1 & 2 \\ 3 & 4 \end{pmatrix}$. Its characteristic polynomial is $p(\lambda) = \lambda^2 - 5\lambda - 2$. By Cayley--Hamilton, $A^2 - 5A - 2I = 0$.
\end{example}

\section{Combinatorics and Elementary Number Theory}
\textbf{Prerequisite Knowledge}: Permutations, combinations, primes, modular arithmetic.

\begin{definition}[Generating Function]
The ordinary generating function for a sequence $(a_0, a_1, a_2, \dots)$ is the formal power series $A(x) = \sum_{n=0}^{\infty} a_n x^n$.
\end{definition}

\begin{definition}[Primitive Root]
An integer $g$ is a primitive root modulo $n$ if the multiplicative order of $g$ modulo $n$ is equal to $\varphi(n)$ (Euler's totient function). This means the powers of $g$ generate the group $(\mathbb{Z}/n\mathbb{Z})^\times$.
\end{definition}

\begin{theorem}[Chinese Remainder Theorem]
Let $m_1, m_2, \dots, m_k$ be pairwise coprime positive integers. For any integers $a_1, a_2, \dots, a_k$, the system of simultaneous congruences $x \equiv a_i \pmod{m_i}$ for $1 \le i \leq k$ has a unique solution modulo $M = m_1 m_2 \dots m_k$.
\end{theorem}

\begin{example}
Find $x$ such that $x \equiv 2 \pmod 3$ and $x \equiv 3 \pmod 5$. Since 3 and 5 are coprime, the CRT guarantees a unique solution modulo 15. The solution is $x = 8$.
\end{example}